\section{Extended Attack-Resilience Results}
\label{app:attacks}

Table~\ref{tab:attack-full} reports the full ten-scenario single-actor
attack suite. A1 and A2 ran on the earlier four-validator cluster ($n{=}4$);
all other scenarios ran on the seven-validator testbed ($n{=}7$). The column
\emph{Byz.\ VP} gives the share of voting power held by the misbehaving
validator: one of seven validators is 14.3\%, and the attacker containers of
A1 and A2 are non-validators (0\%).

\begin{table}[h]
\caption{Full Attack-Resilience Suite (live)}
\label{tab:attack-full}
\centering
\footnotesize
\begin{tabular}{@{}p{0.04\linewidth}p{0.18\linewidth}p{0.03\linewidth}p{0.07\linewidth}p{0.06\linewidth}p{0.06\linewidth}p{0.36\linewidth}@{}}
\toprule
\textbf{\#} & \textbf{Attack} & \textbf{$n$} & \textbf{Byz.\ VP} & \textbf{Safety} & \textbf{Diverg.} & \textbf{Attack-specific signal} \\
\midrule
A1 & Eclipse & 4 & 0\% & True & 0 & Filter holds 8/8, nil-prevote ratio 0.0 throughout \\
A2 & Simulated multi-subnet Sybil & 4 & 0\% & True & 0 & Filter holds 8/8 despite 12 attackers \\
A3 & Data withholding & 7 & 14.3\% & True & 0 & --- \\
A4 & Forged Bitcoin receipt & 7 & 14.3\% & True & 0 & --- \\
A5 & Withdrawal during \textsc{Sovereign} & 7 & --- & True & 0 & Withdrawal correctly blocked; transaction never commits \\
A6 & Malicious proposer & 7 & 14.3\% & True & 0 & Forged FSM state rejected by all honest validators (Section~\ref{subsec:rq1}) \\
A7 & Censorship & 7 & 14.3\% & True & 0 & Censoring confirmed real (target tx omitted 9/9 proposer turns); honest reject signals stayed at 0 since the next honest proposer included the tx before the ignore-round threshold fired \\
A8 & Combined (A6 + A1) & 7 & 14.3\% & True & 0 & Safety held under simultaneous overlap \\
A9 & Double-signing & 7 & 14.3\% & True & 0 & 6/6 honest validators detected real duplicate-vote evidence; 1-block detection latency \\
A10 & Timeout flood & 7 & 14.3\% & True & 0 & Cadence held (0.770$\to$0.706 blocks/s); rate limiter dropped 0/0/0/0/0/0 at moderate (20 msg/s) intensity \\
\bottomrule
\end{tabular}
\end{table}

A3 and A4 were instrumented only for safety and divergence, without an
attack-specific metric such as A5's blocked-withdrawal flag or A9's
detection latency.

At 25$\times$ its design threshold (500 msg/s vs.\ a 20 msg/s per-peer cap),
the A10 timeout-flood attack still held cadence (0.825$\to$0.789 blocks/s)
while the rate limiter rejected the large majority of flood traffic
(28{,}071 of $\sim$30{,}000 attempted messages on one validator); validator
CPU usage roughly tripled (2.4--3.3\%$\to$7.85--9.51\%, peak 22.40\%) but
stayed well short of saturation, and memory usage stayed flat. No attack in
Table~\ref{tab:attack-full} drove block production toward zero, and resource
usage under the P2P swarm attacks (A1/A2) stayed modest (under 36\% CPU,
under 147\,MiB memory per validator).

\FloatBarrier
